\subsection*{2.3 Measure Functions}

We will show in the last section of this chapter that it can lead to logical contradictions if we attempt to define a probability measure for all subsets in a sample space. To avoid this, probability theory adopts the solution of defining the probability only for some selected subsets.

The $\sigma$-algebra is significant because it serves as the domain of a measure function in probability theory. The measure function is a set function that takes a set as input and returns a nonnegative real number as output, satisfying some axioms.

\begin{defbox}{2.5}
Let $\mathcal{F}$ be a $\sigma$-field on a sample space $\Omega$. A set function $m$ from $\mathcal{F}$ to $[0,\infty]$ is called a \textit{measure} if it satisfies the following properties:
\begin{enumerate}[label=\arabic*.]
    \item $m(\emptyset)=0$.
    \item For any sequence of mutually disjoint $A_i\in \mathcal{F}$, for $i=1,2,3,\dots$, we have
    \[
    m\left(\biguplus_{i=1}^{\infty} A_i\right)
    =
    \sum_{i=1}^{\infty} m(A_i).
    \]
\end{enumerate}

This property is known as the $\sigma$-\textit{additive} property.

A measure $m$ is said to be \textit{finite} if $m(\Omega)$ is finite. If $m(\Omega)=1$, then $m$ is called a \textit{probability measure}.

A \textit{measure space} is a triple $(\Omega,\mathcal{F},m)$ where $\mathcal{F}$ is a $\sigma$-field on $\Omega$ and $m$ is a measure on $\mathcal{F}$. A \textit{probability space} is a measure space $(\Omega,\mathcal{F},m)$ when $\mathcal{F}$ is a $\sigma$-field on $\Omega$ and $m$ is a probability measure.
\end{defbox}

\textbf{Convention with $\infty$} \\
In Definition 2.5, the measure function $m$ may take the value $\infty$. We adopt the usual convention for arithmetic operations involving of $\infty$, namely:
\begin{itemize}
    \item $c+\infty=\infty$ for all real numbers $c$.
    \item $\infty+\infty=\infty$.
\end{itemize}

\textbf{Example 2.3.1 (Counting Measure)} \\
Let $\Omega$ be any countable set, and let the power set of $\Omega$, $\mathcal{P}(\Omega)$, be the $\sigma$-field $\mathcal{F}$. Define the \textit{counting measure} $\mu$ of a set $A\subseteq \Omega$ to be the cardinality of $A$. Then $(\Omega,\mathcal{P}(\Omega),\mu)$ is a measure space.

\textbf{Example 2.3.2 (Discrete Probability Space)} \\
Let $\Omega=\mathbb{N}$. Suppose $p_1,p_2,p_3,\dots$ is a sequence of nonnegative real numbers such that $\sum_{i=1}^{\infty} p_i=1$. For any $E\subseteq \Omega$, let $P(E)=\sum_{i\in E} p_i$. Then $(\Omega,\mathcal{P}(\Omega),P)$ is a probability space.

The first two properties below are elementary, and they hold more generally for any pre-measure (which will be defined in the next chapter).

\begin{thmbox}{2.2 (Monotonicity)}
Let $(\Omega,\mathcal{F},\mu)$ be a measure space, and suppose $A$ and $B$ are $\mathcal{F}$-measurable sets such that $A\subseteq B$. Then $\mu(A)\le \mu(B)$.
\end{thmbox}

\textit{Proof} Since $A$ and $B\setminus A$ are disjoint and $B=A\cup (B\setminus A)$, we have
\[
\mu(B)=\mu(A)+\mu(B\setminus A).
\]
Since $\mu(B\setminus A)$ is nonnegative, it follows that $\mu(B)\ge \mu(A)$. \hfill $\square$

\begin{thmbox}{2.3 (Finite Subadditivity)}
Let $(\Omega,\mathcal{F},\mu)$ be a measure space, and suppose $A$ and $B$ are $\mathcal{F}$-measurable sets. Then
\[
\mu(A\cup B)\le \mu(A)+\mu(B).
\]
\end{thmbox}

\textit{Proof} $\mu(A\cup B)=\mu(A \uplus (B\setminus A))=\mu(A)+\mu(B\setminus A)\le \mu(A)+\mu(B)$. \hfill $\square$

The next general property is about an increasing and decreasing sequence of events.

\begin{defbox}{2.6}
A sequence of sets $A_i$, for $i=1,2,3,\dots$, is said to be \textit{increasing} if $A_1\subseteq A_2\subseteq A_3\subseteq \cdots$, or \textit{decreasing} if $A_1\supseteq A_2\supseteq A_3\supseteq \cdots$.
\end{defbox}

\begin{thmbox}{2.4 (Lower Semi-Continuity)}
Suppose $A_1\subseteq A_2\subseteq A_3\subseteq \cdots$ is a sequence of increasing sets in $\mathcal{F}$ and $(\Omega,\mathcal{F},\mu)$ is a measure space. Then
\[
\lim_{k\to\infty} \mu(A_k)=\mu\left(\bigcup_{i=1}^{\infty} A_i\right).
\]
\end{thmbox}

\textit{Proof} Let $B_1=A_1$ and $B_i=A_i\setminus A_{i-1}$ for $i\ge 2$. The sets $B_i$'s are mutually disjoint by construction. Moreover, we have $\biguplus_{i=1}^{\infty} B_i=\cup_{i=1}^{\infty} A_i$. This gives
\[
\lim_{k\to\infty} \mu(A_k)
=
\lim_{k\to\infty} \mu(B_1 \uplus B_2 \uplus \cdots \uplus B_k)
\]
\[
=
\lim_{k\to\infty} \sum_{i=1}^{k} \mu(B_i)
\triangleq
\sum_{i=1}^{\infty} \mu(B_i)
=
\mu\left(\biguplus_{i=1}^{\infty} B_i\right)
=
\mu\left(\bigcup_{i=1}^{\infty} A_i\right),
\]
where the second last equality follows from the countable additivity of measure $\mu$. \hfill $\square$

\begin{thmbox}{2.5 (Upper Semi-Continuity)}
Suppose $A_1\supseteq A_2\supseteq A_3\supseteq \cdots$ is a sequence of decreasing sets in $\mathcal{F}$, and $(\Omega,\mathcal{F},\mu)$ is a measure space. Furthermore, suppose $\mu(A_1)<\infty$. Then
\[
\lim_{k\to\infty} \mu(A_k)=\mu\left(\bigcap_{i=1}^{\infty} A_i\right).
\]
\end{thmbox}

\textit{Proof} Apply Theorem 2.4 to the sequence of events $E_i\triangleq A_1\setminus A_i$, for $i=1,2,3,\dots$, and exploit that fact that $\mu(E_i)=\mu(A_1)-\mu(A_i)$. \hfill $\square$

The properties in Theorems 2.4 and 2.5 are also known as \textit{continuity from below} and \textit{continuity from above}, respectively.

It is a customary to use the notation $A_k \nearrow A$ to signify that $(A_k)_{k=1}^{\infty}$ is a sequence of increasing sets with union $A$. Similarly, we write $A_k \searrow A$ for a sequence of decreasing sets with intersection $A$.
